%
% This work is licensed under a Creative Commons Attribution-ShareAlike 4.0 International License.
%

\documentclass[10pt, letterpaper, twoside, openany]{article}

\title{Introduction to Quantum Mechanics -- Worksheet 4}
\author{Joseph D. MacMillan}
\date{}


\usepackage{graphicx}
\usepackage{color} 
\usepackage{amsmath, amssymb}
\usepackage{libertinus}
\usepackage{microtype}
\usepackage{layout}
\usepackage[most]{tcolorbox}
\tcbuselibrary{skins,breakable}
\usepackage[hang, small, bf]{caption}
\captionsetup[table]{position=top}
\captionsetup[figure]{position=bottom}
\usepackage[hyphens]{url}
\usepackage[hidelinks]{hyperref}
\hypersetup{pdftitle={Introduction to Quantum Mechanics -- Worksheet 4}}
\hypersetup{pdfauthor={Joseph D. MacMillan}}
\urlstyle{rm}
\usepackage[shortlabels]{enumitem}
\usepackage[margin=1in]{geometry}
\usepackage{fancyhdr}




\newcommand{\uu}{\symbfup{u}}
\newcommand{\grad}{\symbfup{\nabla}}
\newcommand{\curl}{\symbfup{\nabla} \times}
\newcommand{\dfdx}[2]{\frac{\partial {#1}}{\partial {#2}}}
\newcommand{\ddfdx}[2]{\frac{\partial^2 {#1}}{\partial {#2}^2}}
\newcommand{\vort}{\symbfup{\omega}}
\renewcommand\vec{\symbfup}
\newcommand{\unit}[1]{\hat{\vec{#1}}}
\newcommand{\U}{\mathbb{U}}
\newcommand{\V}{\mathbb{V}}
\newcommand{\LL}{\mathbb{L}^2}
\newcommand{\LLLL}{\mathbb{L}^4}

\newcommand{\ket}[1]{| #1 \rangle}
\newcommand{\bra}[1]{\langle #1 |}
\newcommand{\braket}[2]{\langle #1 | #2 \rangle}
\newcommand{\avg}[1]{\langle #1 \rangle}


%\DeclareMathAlphabet{\mathbfsf}{\encodingdefault}{\sfdefault}{bx}{sl}
\newcommand{\tens}[1]{\mathbfsfit{#1}}


\newcounter{example}[section]
\def\theexample{\thechapter.\arabic{example}}
\newenvironment{example}[1][ ]{\refstepcounter{example}
\begin{tcolorbox}[breakable, sharp corners, boxrule = 0pt, frame empty, opacityframe=0, parbox=false]
\textbf{Example \theexample \ -- #1.}
}
{ 
\end{tcolorbox}
}

\newenvironment{theorem}[1][ ]{
\begin{tcolorbox}[breakable, sharp corners, boxrule = 0pt, frame empty, opacityframe=0, parbox=false]
\textbf{#1.}
}
{ 
\end{tcolorbox}
}

\newcounter{problem}[section]
\def\theproblem{\thechapter.\arabic{problem}}
\newenvironment{problem}[1][ ]{\refstepcounter{problem} \noindent \textbf{Problem \theproblem \ -- #1.}}{\vspace{0.1in}}

\renewcommand{\arraystretch}{2.5}

\begin{document}

\pagestyle{fancy}
\fancyhead[L]{\emph{Introduction to Quantum Mechanics}}
\fancyhead[R]{Worksheet 4} 

\Large \noindent Worksheet 4

\noindent \LARGE \textbf{One Big Example}

\normalsize

\vspace{0.2in}

\noindent Quantum States  \bullet \  Probabilities   \bullet \  Measurement   

\vspace{0.1in}

\noindent \rule{\textwidth}{0.5pt}

\vspace{0.2in}

As a quantum physicist, your job is to prepare an ensemble of spin-1/2 particles, all in the same state given by
\[
\ket{\psi} = \ket{+} - 2i \ket{-}.
\]
You then make a series of measurements on the ensemble, each time returning it to this original state.

\begin{enumerate}
\item Would you say that the particle in your ensemble are more ``spin-up'' or more ``spin-down''?  Why?

\vspace{1in}

\item Normalize the state.  Why do we need to work with states that are normalized?

\vspace{2in}

\item Write the state as a bra, and write both the ket and the bra in matrix notation.

\newpage

\item While a particle is in the state $\ket{\psi}$, what is its spin component along the $z$ axis, $S_z$?  

\vspace{1.5in}

\item What are the possibilities and probabilities when measuring 
\begin{enumerate}
\item $S_z$?

\vspace{1.5in}

\item $S_x$?

\vspace{2.5in}

\item $S_y$?


\end{enumerate}

\newpage

\item If a measurement of $S^2$ is made, what will be measured and with what probability?

\vspace{2in}

\item Calculate the expectation values
\begin{enumerate}
\item $\avg{S_z}$ and $\avg{S_z^2}$.

\vspace{3in}

\item $\avg{S_x}$ and $\avg{S_x^2}$.

\newpage

\item $\avg{S_y}$ and $\avg{S_y^2}$.
\end{enumerate}

\vspace{3in}

\item Calculate $\sigma_{S_x}$, $\sigma_{S_y}$, and $\sigma_{S_z}$, and show that the uncertainty principle 
\[
\sigma_{S_x} \sigma_{S_y} \ge \frac{\hbar}{2} | \avg{S_z} |
\]
holds for this state.

\end{enumerate}

\end{document}
